% Lean: EconLean/Macro/Labor/HosiosDiscounted.lean
\chapter{Matching with Discounting: the Hosios Condition for $r > 0$}
\label{chap:hosios_discounted}

Sources: Pissarides, \emph{Equilibrium Unemployment Theory} (2000), \S8.1; Hosios (1990); Ljungqvist--Sargent
ch.~28. Chapter~\ref{chap:labor} proves Hosios for the steady-state ($r \to 0$) planner. Here time is discrete,
payoffs are discounted by $\delta = 1/(1+r) \in (0,1)$, and the planner maximizes the present value of net output
along the whole transition path. The matching primitives $A, \eta, s, c, y, b$, the rates
$f(\theta) = A\theta^{1-\eta}$, $q(\theta) = A\theta^{-\eta}$ and the job-creation function $jc$ are those of
Section~\ref{sec:dmp}.

\begin{definition}[Discrete-time matching economy]
\label{def:hd_economy}
\lean{Econ.DMP.Discrete, Econ.DMP.Discrete.rate}
\leanok
\uses{def:dmp_economy}
A DMP economy together with $\delta \in (0,1)$ and $s \le 1$. The implied interest rate is $r = 1/\delta - 1$.
\end{definition}

\begin{lemma}[Elementary facts]
\label{lem:hd_basic}
\lean{Econ.DMP.Discrete.rate_pos, Econ.DMP.Discrete.findRate_nonneg}
\leanok
\uses{def:hd_economy}
$r > 0$, and $f(\theta) \ge 0$ for $\theta \ge 0$.
\end{lemma}
\begin{proof}
\leanok
$1/\delta > 1$ since $0 < \delta < 1$; $f(\theta) = A\theta^{1-\eta}$ with $A > 0$ and $\theta^{1-\eta} \ge 0$.
\end{proof}

\begin{definition}[Stationary search equilibrium]
\label{def:hd_searchEqm}
\lean{Econ.DMP.Discrete.IsSearchEqm}
\leanok
\uses{def:hd_economy}
$(\theta, J, E, U, w)$ with $\theta > 0$, matching rates that are probabilities ($f(\theta) \le 1$,
$q(\theta) \le 1$), and
$J = y - w + \delta(1-s)J$, $E = w + \delta((1-s)E + sU)$, $U = b + \delta(f(\theta)E + (1 - f(\theta))U)$,
Nash sharing $(1-\beta)(E - U) = \beta J$, and free entry $c = \delta q(\theta) J$.
\end{definition}

\begin{lemma}[Job creation at $r = 1/\delta - 1$]
\label{lem:hd_jc_rate}
\lean{Econ.DMP.Discrete.jc_rate}
\leanok
\uses{def:hd_economy, def:dmp_jc}
For $\theta > 0$: $jc(r, \beta, \theta)\cdot \delta q(\theta) = c\,(1 - \delta(1-s) + \delta\beta f(\theta))$.
\end{lemma}
\begin{proof}
\leanok
Write $A = q(\theta)\theta^\eta$ and $f(\theta) = \theta q(\theta)$ (Lemma \texttt{findRate\_eq}); substitute into
$jc = c(r+s)\theta^\eta/A + \beta c\theta$ with $r + s = (1 - \delta(1-s))/\delta$ and clear denominators.
\end{proof}

\begin{theorem}[Equilibrium is job creation at $r = 1/\delta - 1$]
\label{thm:hd_isSearchEqm_iff}
\lean{Econ.DMP.Discrete.isSearchEqm_iff, Econ.DMP.Discrete.existsUnique_searchEqm}
\leanok
\uses{def:hd_searchEqm, lem:hd_jc_rate, thm:dmp_existsUnique}
For $\beta \ge 0$: a search equilibrium exists at $\theta$ iff $\theta$ is a DMP equilibrium with interest rate
$r = 1/\delta - 1$ and $f(\theta), q(\theta) \le 1$. Hence for $\beta \in [0,1)$, if the job-creation solution has
$f, q \le 1$, it is the unique search equilibrium tightness.
\end{theorem}
\begin{proof}
\leanok
($\Rightarrow$) Let $S = J + (E - U)$ be the match surplus. Subtracting the $U$ equation from the $E$ equation
gives $(E-U)(1 - \delta(1 - s - f)) = w - b$; the firm equation gives $(1 - \delta(1-s))J = y - w$; Nash sharing
gives $E - U = \beta S$. One linear combination of these four equations gives
$J(1 - \delta(1-s) + \delta\beta f) = (1-\beta)(y-b)$. Multiplying by $\delta q$ and using free entry
$c = \delta q J$ gives $c(1 - \delta(1-s) + \delta\beta f) = \delta q (1-\beta)(y-b)$, which by
Lemma~\ref{lem:hd_jc_rate} is $jc = (1-\beta)(y-b)$ after cancelling $\delta q > 0$.
($\Leftarrow$) Put $K = 1 - \delta(1-s) + \delta\beta f > 0$, $S = (y-b)/K$, $J = (1-\beta)S$, $E - U = \beta S$,
$U = (b + \delta f\beta S)/(1-\delta)$, $E = U + \beta S$, $w = y - (1 - \delta(1-s))J$, and check each equation;
free entry follows from Lemma~\ref{lem:hd_jc_rate}. The probability bounds pass through unchanged in both
directions. Uniqueness then follows from Theorem~\ref{thm:dmp_existsUnique} with $r > 0$.
\end{proof}

\begin{definition}[The planner's problem]
\label{def:hd_planner}
\lean{Econ.DMP.Discrete.next, Econ.DMP.Discrete.upath, Econ.DMP.Discrete.reward, Econ.DMP.Discrete.Feasible, Econ.DMP.Discrete.welfare, Econ.DMP.Discrete.IsOptimal}
\leanok
\uses{def:hd_economy}
Unemployment moves by $u' = u + s(1-u) - f(\theta)u$; period net output is $y(1-u) + bu - c\theta u$. A tightness
path is feasible if $\theta_t \ge 0$, $f(\theta_t) \le 1$ and $q(\theta_t) \le 1$ (matches exceed neither the
unemployed nor vacancies; at $\theta = 0$ Lean's $q(0) = 0$ makes the last condition vacuous). Welfare from $u_0$ is
$\sum_t \delta^t\,\mathrm{reward}(u_t, \theta_t)$, and a path is optimal if it is feasible and no feasible path
gives more welfare.
\end{definition}

\begin{lemma}[{Feasible paths stay in $[0,1]$, are bounded, and have summable welfare}]
\label{lem:hd_feasible}
\lean{Econ.DMP.Discrete.upath_mem, Econ.DMP.Discrete.le_bound, Econ.DMP.Discrete.summable_welfare}
\leanok
\uses{def:hd_planner, lem:hd_basic}
If $u_0 \in [0,1]$ and the path is feasible, then $u_t \in [0,1]$, $\theta_t \le (1/A)^{1/(1-\eta)}$, and the
welfare series converges absolutely.
\end{lemma}
\begin{proof}
\leanok
$u' = s(1-u) + (1-f)u$ is a convex combination of $s \in (0,1]$ and $1 - f \in [0,1]$, by induction. From
$A\theta^{1-\eta} \le 1$, raise to the power $1/(1-\eta)$ (\texttt{rpow\_le\_rpow}, \texttt{rpow\_rpow\_inv}).
Then $|\mathrm{reward}| \le |y| + |b| + c(1/A)^{1/(1-\eta)}$, and the series is dominated by a geometric series
(\texttt{Summable.of\_norm\_bounded}).
\end{proof}

\begin{definition}[Shadow value, value function, one-period gap]
\label{def:hd_value}
\lean{Econ.DMP.Discrete.shadow, Econ.DMP.Discrete.value, Econ.DMP.Discrete.gap}
\leanok
\uses{def:hd_planner}
Given $\theta_p$, let $\lambda = (y-b)/(1 - \delta(1-s) + \delta\eta f(\theta_p))$, $V(u) = \alpha - \lambda u$ with
$\alpha = (y - \delta\lambda s)/(1-\delta)$, and $\mathrm{gap}(u, \theta) = \mathrm{reward}(u,\theta) + \delta V(u') - V(u)$.
\end{definition}

\begin{lemma}[The gap]
\label{lem:hd_gap}
\lean{Econ.DMP.Discrete.gap_eq, Econ.DMP.Discrete.findRate_tangent, Econ.DMP.Discrete.shadow_pos, Econ.DMP.Discrete.foc, Econ.DMP.Discrete.gap_le, Econ.DMP.Discrete.gap_lt, Econ.DMP.Discrete.gap_self}
\leanok
\uses{def:hd_value, lem:hd_jc_rate}
Let $\theta_p$ solve job creation at $r$ with $\beta = \eta$. Then
$\mathrm{gap}(u,\theta) = u\,[(\delta\lambda f(\theta) - c\theta) - (\delta\lambda f(\theta_p) - c\theta_p)]$, $\lambda > 0$,
and $\delta\lambda(1-\eta)q(\theta_p) = c$. Consequently $\mathrm{gap}(u, \theta) \le 0$ for $u, \theta \ge 0$, with strict
inequality if $u > 0$ and $\theta \ne \theta_p$, and $\mathrm{gap}(u, \theta_p) = 0$.
\end{lemma}
\begin{proof}
\leanok
From Lemma~\ref{lem:hd_jc_rate} with $\beta = \eta$, $c\,K = \delta q(1-\eta)(y-b) = \delta q (1-\eta)\lambda K$, so
$c = \delta\lambda(1-\eta)q(\theta_p)$ and $c\theta_p = \delta\lambda(1-\eta)f(\theta_p)$. Substituting these and the
definitions of $\lambda$ and $\alpha$ makes the constant term vanish and leaves the stated multiple of $u$
(\texttt{linear\_combination}). Concavity of $f$: with $x = \theta/\theta_p - 1 \ge -1$,
Bernoulli's inequality $(1+x)^{1-\eta} \le 1 + (1-\eta)x$ (\texttt{rpow\_one\_add\_le\_one\_add\_mul\_self},
strict version for $x \ne 0$) gives $f(\theta) \le f(\theta_p) + (1-\eta)q(\theta_p)(\theta - \theta_p)$. Multiply by
$\delta\lambda > 0$ and use the first-order condition.
\end{proof}

\begin{lemma}[Telescoping and limits]
\label{lem:hd_telescope}
\lean{Econ.DMP.Discrete.telescope, Econ.DMP.Discrete.tendsto_tail, Econ.DMP.Discrete.tendsto_gaps}
\leanok
\uses{def:hd_value, lem:hd_feasible}
$\sum_{t<T}\delta^t \mathrm{reward}_t + \delta^T V(u_T) = V(u_0) + \sum_{t<T}\delta^t\mathrm{gap}_t$; along a feasible path
$\delta^T V(u_T) \to 0$; hence $V(u_0) + \sum_{t<T}\delta^t \mathrm{gap}_t \to \mathrm{welfare}$.
\end{lemma}
\begin{proof}
\leanok
Induction on $T$. $V(u_T)$ is bounded since $u_T \in [0,1]$, and $\delta^T \to 0$ (squeeze). Add the two limits.
\end{proof}

\begin{theorem}[The efficient path]
\label{thm:hd_optimal}
\lean{Econ.DMP.Discrete.welfare_stationary, Econ.DMP.Discrete.welfare_le, Econ.DMP.Discrete.welfare_lt_value, Econ.DMP.Discrete.isOptimal_stationary, Econ.DMP.Discrete.welfare_lt}
\leanok
\uses{lem:hd_gap, lem:hd_telescope}
Let $\theta_p$ solve job creation at $r = 1/\delta - 1$ with $\beta = \eta$, $f(\theta_p) \le 1$ and
$q(\theta_p) \le 1$. From every
$u_0 \in [0,1]$ the stationary path $\theta_t = \theta_p$ is optimal, with welfare $V(u_0)$. If $u_0 > 0$, every
feasible path with $\theta_0 \ne \theta_p$ gives strictly less.
\end{theorem}
\begin{proof}
\leanok
For the stationary path every gap is $0$, so welfare is $V(u_0)$ by Lemma~\ref{lem:hd_telescope} and uniqueness of
limits. For any feasible path all gaps are $\le 0$, so welfare $\le V(u_0)$ (\texttt{le\_of\_tendsto\_of\_tendsto'}).
If $\theta_0 \ne \theta_p$ and $u_0 > 0$, then for $T \ge 1$ the partial sums of gaps are at most
$\mathrm{gap}_0 < 0$, so welfare $\le V(u_0) + \mathrm{gap}_0 < V(u_0)$.
\end{proof}

\begin{theorem}[Hosios with discounting]
\label{thm:hd_hosios}
\lean{Econ.DMP.Discrete.hosios}
\leanok
\uses{thm:hd_optimal, thm:hd_isSearchEqm_iff, thm:dmp_strictAnti_beta}
Let $\beta \in [0,1)$, $\theta_e$ a search equilibrium tightness, $\theta_p$ the job-creation solution with
$\beta = \eta$ and $f(\theta_p), q(\theta_p) \le 1$, and $u_0 \in (0, 1]$. The stationary equilibrium path is optimal iff
$\beta = \eta$.
\end{theorem}
\begin{proof}
\leanok
By Theorem~\ref{thm:hd_isSearchEqm_iff}, $\theta_e$ solves job creation at $r$. ($\Rightarrow$) If $\theta_e \ne \theta_p$,
Theorem~\ref{thm:hd_optimal} gives welfare$(\theta_e) <$ welfare$(\theta_p)$, contradicting optimality; so
$\theta_e = \theta_p$, and if $\beta \ne \eta$ Theorem~\ref{thm:dmp_strictAnti_beta} makes the two equilibria differ.
($\Leftarrow$) With $\beta = \eta$, uniqueness gives $\theta_e = \theta_p$, which is optimal.
\end{proof}

\begin{remark}
In discrete time $f(\theta)$ and $q(\theta)$ are probabilities, so the planner is restricted to
$f(\theta_t), q(\theta_t) \le 1$, and the theorem assumes both at $\theta_p$; the Lean file gives a calibration ($A = 1/4$, $\eta = 1/2$, $s = 1/10$,
$c = y = 1$, $b = 0$, $\delta = 9/10$) where it holds.
\end{remark}
